\uuid{eeiI}
\titre{Changements de variables}
\auteur{hanani ; mzari}
\organisation{exo7}
\licence{CC BY-NC-SA 4.0}
\datecreate{2020-01-20}
\niveau{L1}
\module{Analyse}
\chapitre{Calcul d'intégrales}
\sousChapitre{Autre}
\difficulte{3}
\theme{Calculs d'intégrales, changements de variables}
\exo7qcmid{lille-2.integrales.025}

\contenu{
\begin{question}
\qclassification{127.99}
\qtags{motcle=changements de variables}
Parmi les affirmations suivantes, cocher celles qui sont vraies :
\begin{answers}
\good{La dérivée de $\tan x$ sur $\displaystyle ]-\pi/2,\pi/2[$ est $1+\tan ^2x$.}
\bad{Le changement de variable $t=\tan x$ donne $\displaystyle \int _0^{\pi/3}\frac{\mathrm{d}x}{1+2\cos ^2x}=\int _0^{\pi/3}\frac{\mathrm{d}t}{3+t^2}$.}
\bad{Une primitive de $\displaystyle \frac{1}{3+t^2}$ sur $\Rr$ est $\displaystyle \frac{1}{3}\arctan \left(\frac{t}{\sqrt{3}}\right)$.}
\good{$\displaystyle \int _0^{\pi/3}\frac{\mathrm{d}x}{1+2\cos ^2x}=\frac{\pi}{4\sqrt{3}}$.}
\end{answers}
\vskip2mm
\begin{explanations}
Avec $t=\tan x\Rightarrow \mathrm{d}t=(1+\tan ^2x)\, \mathrm{d}x$, on a : $t(0)=0$, $t(\pi/3)=\sqrt{3}$ et
$$\int _0^{\pi/3}\frac{\mathrm{d}x}{1+2\cos ^2x}=\int _0^{\sqrt{3}}\frac{\mathrm{d}t}{3+t^2}=\frac{1}{\sqrt{3}}\left[\arctan \left(\frac{t}{\sqrt{3}}\right)\right]_0^{\sqrt{3}}=\frac{\pi}{4\sqrt{3}}.$$
\end{explanations}
\end{question}
}
